% GENERATED FILE. Edit canonical lessons, then run npm run generate:books.
% canonical-source-hash: fnv1a64:04df84703daaba3b
% canonical-lessons: KA-C303-jagali, KA-C303-atta, KA-C303-cilaka, KA-C303-hidike, KA-C303-vasane

\chapter{Raised veranda, Loft, Door latch, Handle, Smell}
\label{ch:ka-raised-veranda-loft-door-latch-handle-smell}
\hlchaptermodality{303}

\begin{chapteropening}
\textbf{By the end of this chapter:} \emph{I can say a raised veranda, a loft, a door latch, a handle and a smell.}

Complete the last lesson of chapter 303: I can say a raised veranda, a loft, a door latch, a handle and a smell.
\end{chapteropening}

\section[\texorpdfstring{\kn{ಜಗಲಿ} (jagali)}{ಜಗಲಿ (jagali)}]{\kn{ಜಗಲಿ} (jagali) — a raised veranda, a front porch}

\label{lesson:KA-C303-jagali}



\begin{quote}
Before the new one: say the Kannada for cement, then the Kannada for bamboo.
\end{quote}

\subsection*{You'll want to know: \kn{ಜಗಲಿ}}
\textbf{\kn{ಜಗಲಿ}} — \emph{jagali} — \textquotedblleft{}a raised veranda, a front porch\textquotedblright{}.

\begin{grammarlens}[title={What you've built}]
One more word for this chapter.
\end{grammarlens}

\subsection*{Your turn}
\begin{itemize}
\raggedright
  \item \emph{Say it:} \emph{jagali}
  \item \emph{Say it:} \emph{jagali}, once more
  \item \emph{Say it:} say \emph{jagali}
  \item \emph{Recall:} say the Kannada for cement, then the Kannada for bamboo, then say \emph{jagali} again
\end{itemize}

\begin{tcolorbox}[breakable,colback=teal!4,colframe=teal!35!black,title={Before you move on}]
What does \kn{ಜಗಲಿ} mean? (\textquotedblleft{}A raised veranda, a front porch\textquotedblright{}.) Say it once more.
\end{tcolorbox}
\section[\texorpdfstring{\kn{ಅಟ್ಟ} (aṭṭa)}{ಅಟ್ಟ (aṭṭa)}]{\kn{ಅಟ್ಟ} (aṭṭa) — a loft, an attic}

\label{lesson:KA-C303-atta}



\begin{quote}
Before the new one: say the Kannada for bamboo, then the Kannada for a raised veranda.
\end{quote}

\subsection*{You'll want to know: \kn{ಅಟ್ಟ}}
\textbf{\kn{ಅಟ್ಟ}} — \emph{aṭṭa} — \textquotedblleft{}a loft, an attic\textquotedblright{}.

\begin{grammarlens}[title={What you've built}]
One more word for this chapter.
\end{grammarlens}

\subsection*{Your turn}
\begin{itemize}
\raggedright
  \item \emph{Say it:} \emph{aṭṭa}
  \item \emph{Say it:} \emph{aṭṭa}, once more
  \item \emph{Say it:} say \emph{aṭṭa}
  \item \emph{Recall:} say the Kannada for bamboo, then the Kannada for a raised veranda, then say \emph{aṭṭa} again
\end{itemize}

\begin{tcolorbox}[breakable,colback=teal!4,colframe=teal!35!black,title={Before you move on}]
What does \kn{ಅಟ್ಟ} mean? (\textquotedblleft{}A loft, an attic\textquotedblright{}.) Say it once more.
\end{tcolorbox}
\section[\texorpdfstring{\kn{ಚಿಲಕ} (cilaka)}{ಚಿಲಕ (cilaka)}]{\kn{ಚಿಲಕ} (cilaka) — a door latch, a bolt}

\label{lesson:KA-C303-cilaka}



\begin{quote}
Before the new one: say the Kannada for a raised veranda, then the Kannada for a loft.
\end{quote}

\subsection*{You'll want to know: \kn{ಚಿಲಕ}}
\textbf{\kn{ಚಿಲಕ}} — \emph{cilaka} — \textquotedblleft{}a door latch, a bolt\textquotedblright{}.

\begin{grammarlens}[title={What you've built}]
One more word for this chapter.
\end{grammarlens}

\subsection*{Your turn}
\begin{itemize}
\raggedright
  \item \emph{Say it:} \emph{cilaka}
  \item \emph{Say it:} \emph{cilaka}, once more
  \item \emph{Say it:} say \emph{cilaka}
  \item \emph{Recall:} say the Kannada for a raised veranda, then the Kannada for a loft, then say \emph{cilaka} again
\end{itemize}

\begin{tcolorbox}[breakable,colback=teal!4,colframe=teal!35!black,title={Before you move on}]
What does \kn{ಚಿಲಕ} mean? (\textquotedblleft{}A door latch, a bolt\textquotedblright{}.) Say it once more.
\end{tcolorbox}
\section[\texorpdfstring{\kn{ಹಿಡಿಕೆ} (hiḍike)}{ಹಿಡಿಕೆ (hiḍike)}]{\kn{ಹಿಡಿಕೆ} (hiḍike) — a handle}

\label{lesson:KA-C303-hidike}



\begin{quote}
Before the new one: say the Kannada for a loft, then the Kannada for a door latch.
\end{quote}

\subsection*{You'll want to know: \kn{ಹಿಡಿಕೆ}}
\textbf{\kn{ಹಿಡಿಕೆ}} — \emph{hiḍike} — \textquotedblleft{}a handle\textquotedblright{}.

\begin{grammarlens}[title={What you've built}]
One more word for this chapter.
\end{grammarlens}

\subsection*{Your turn}
\begin{itemize}
\raggedright
  \item \emph{Say it:} \emph{hiḍike}
  \item \emph{Say it:} \emph{hiḍike}, once more
  \item \emph{Say it:} say \emph{hiḍike}
  \item \emph{Recall:} say the Kannada for a loft, then the Kannada for a door latch, then say \emph{hiḍike} again
\end{itemize}

\begin{tcolorbox}[breakable,colback=teal!4,colframe=teal!35!black,title={Before you move on}]
What does \kn{ಹಿಡಿಕೆ} mean? (\textquotedblleft{}A handle\textquotedblright{}.) Say it once more.
\end{tcolorbox}
\section[\texorpdfstring{\kn{ವಾಸನೆ} (vāsane)}{ವಾಸನೆ (vāsane)}]{\kn{ವಾಸನೆ} (vāsane) — a smell}

\label{lesson:KA-C303-vasane}



\begin{quote}
Before the new one: say the Kannada for a door latch, then the Kannada for a handle.
\end{quote}

\subsection*{You'll want to know: \kn{ವಾಸನೆ}}
\textbf{\kn{ವಾಸನೆ}} — \emph{vāsane} — \textquotedblleft{}a smell\textquotedblright{}.

\begin{grammarlens}[title={What you've built}]
One more word for this chapter.
\end{grammarlens}

\subsection*{Your turn}
\begin{itemize}
\raggedright
  \item \emph{Say it:} \emph{vāsane}
  \item \emph{Say it:} \emph{vāsane}, once more
  \item \emph{Say it:} say \emph{vāsane}
  \item \emph{Recall:} say the Kannada for a door latch, then the Kannada for a handle, then say \emph{vāsane} again
\end{itemize}

\begin{tcolorbox}[breakable,colback=teal!4,colframe=teal!35!black,title={Before you move on}]
What does \kn{ವಾಸನೆ} mean? (\textquotedblleft{}A smell\textquotedblright{}.) Say it once more.
\end{tcolorbox}
